% Ad Familiares 15.21 --- headnote
% ad-familiares-15-21, late 47 BC (probably), Rome

\textit{Cicero to Gaius Trebonius, written from Rome
probably late in 47 BC --- the Perseus dateline reads
\textit{Scr. Romae fort. ex. ann. 707 (47)}, \textit{fort.\
ex.\ ann.}\ marking it as a conjectural late-year placement.
Trebonius, who had served Cicero loyally as quaestor in 60
and broken with the tribune Clodius on his account, was
now on the point of leaving Rome --- the assignment to
Spain glanced at in section 2 is the legateship under
Caesar that would take him to the Further province. He
had sent Cicero a volume gathering Cicero's own witticisms
together with a covering letter; Cicero's reply combines
warm acknowledgement of the gift with a careful defence of
his earlier published assessment of the orator C. Licinius
Calvus, who had recently died and whose letters to Cicero
Trebonius had seen in circulation.}

\textit{The letter is one of Cicero's gentlest performances
in the correspondence with Trebonius, and turns on the
contrast between two kinds of writing: the private letter
to Calvus, written ``no more than this one which you are
now reading, supposing it would ever go abroad,'' and the
book of \textit{bons mots} Trebonius has compiled, which
will. The defence of his praise of Calvus is candid:
genuine respect for a real if limited talent, combined with
the policy of praising in order to spur a young man on. The
catalogue in section 2 of Trebonius's services in 58 ---
the quaestor who took the consuls' part against Clodius
when his own senior colleague would not --- is Cicero's way
of placing the present gift inside a long history of
loyalty, and of registering the personal pang of yet
another friend's departure in a year when the circle of
friends he could speak to freely had grown narrow.}
